% Type A_3 word-level weave Delta Delta -> Delta, drawn in a
% tetrahedral pattern.
% Four-valent commute vertices are suppressed.
\begin{tikzpicture}[
  x=0.94cm,
  y=0.82cm,
  scale=0.80,
  aThreeStrand/.style={line width=0.98pt,line cap=round,line join=round},
  aThreeOne/.style={aThreeStrand,draw=wqteal},
  aThreeTwo/.style={aThreeStrand,draw=wqorange,
    dash pattern=on 3.0pt off 1.45pt},
  aThreeThree/.style={aThreeStrand,draw=wqblue,
    dash pattern=on 0.75pt off 1.35pt},
  aThreeFrame/.style={draw=black!22,line width=0.48pt},
  aThreeT/.style={circle,draw=black!82,fill=blue!12,
    minimum size=2.35mm,inner sep=0pt,line width=0.56pt},
  aThreeH/.style={circle,draw=black!82,fill=white,
    minimum size=3.35mm,inner sep=0pt,line width=0.68pt},
  aThreeVertexLabel/.style={font=\scriptsize,fill=white,
    inner xsep=0.8pt,inner ysep=0.45pt,text=black!88},
  aThreeBoundaryLabel/.style={font=\footnotesize,fill=white,
    inner xsep=1.2pt,inner ysep=0.7pt,text=black!82}
]
\path[use as bounding box] (-6.55,-0.82) rectangle (6.55,10.57);

% The outer triangle has barycentric vertices a,b,c.  The ten weave
% vertices are placed at
%   ((a+(d+1)/3)/3,(b+(d+1)/3)/3,(c+(d+1)/3)/3).
\coordinate (aCorner) at (-6,0);
\coordinate (bCorner) at ( 6,0);
\coordinate (cCorner) at ( 0,{6*sqrt(3)});

\coordinate (t002) at
  (barycentric cs:aCorner=1,bCorner=1,cCorner=7);
\coordinate (t101) at
  (barycentric cs:aCorner=4,bCorner=1,cCorner=4);
\coordinate (t011) at
  (barycentric cs:aCorner=1,bCorner=4,cCorner=4);
\coordinate (t200) at
  (barycentric cs:aCorner=7,bCorner=1,cCorner=1);
\coordinate (t110) at
  (barycentric cs:aCorner=4,bCorner=4,cCorner=1);
\coordinate (t020) at
  (barycentric cs:aCorner=1,bCorner=7,cCorner=1);

\coordinate (h001) at
  (barycentric cs:aCorner=2,bCorner=2,cCorner=5);
\coordinate (h100) at
  (barycentric cs:aCorner=5,bCorner=2,cCorner=2);
\coordinate (h010) at
  (barycentric cs:aCorner=2,bCorner=5,cCorner=2);
\coordinate (h000) at
  (barycentric cs:aCorner=1,bCorner=1,cCorner=1);

% Six boundary points on each side, ordered as 1,2,1,3,2,1.
\coordinate (left1) at
  (barycentric cs:aCorner=1,cCorner=6);
\coordinate (left2) at
  (barycentric cs:aCorner=2,cCorner=5);
\coordinate (left3) at
  (barycentric cs:aCorner=3,cCorner=4);
\coordinate (left4) at
  (barycentric cs:aCorner=4,cCorner=3);
\coordinate (left5) at
  (barycentric cs:aCorner=5,cCorner=2);
\coordinate (left6) at
  (barycentric cs:aCorner=6,cCorner=1);
\coordinate (right1) at
  (barycentric cs:bCorner=1,cCorner=6);
\coordinate (right2) at
  (barycentric cs:bCorner=2,cCorner=5);
\coordinate (right3) at
  (barycentric cs:bCorner=3,cCorner=4);
\coordinate (right4) at
  (barycentric cs:bCorner=4,cCorner=3);
\coordinate (right5) at
  (barycentric cs:bCorner=5,cCorner=2);
\coordinate (right6) at
  (barycentric cs:bCorner=6,cCorner=1);
\coordinate (bottom1) at
  (barycentric cs:aCorner=6,bCorner=1);
\coordinate (bottom2) at
  (barycentric cs:aCorner=5,bCorner=2);
\coordinate (bottom3) at
  (barycentric cs:aCorner=4,bCorner=3);
\coordinate (bottom4) at
  (barycentric cs:aCorner=3,bCorner=4);
\coordinate (bottom5) at
  (barycentric cs:aCorner=2,bCorner=5);
\coordinate (bottom6) at
  (barycentric cs:aCorner=1,bCorner=6);
\draw[aThreeFrame]
  (aCorner)--(cCorner)--(bCorner)--cycle;

% Internal strands of colors 1 and 2.
\draw[aThreeOne] (t002)--(h001);
\draw[aThreeOne] (t101)--(h001);
\draw[aThreeOne] (t011)--(h001);
\draw[aThreeOne] (t101)--(h100);
\draw[aThreeOne] (t200)--(h100);
\draw[aThreeOne] (t110)--(h100);
\draw[aThreeOne] (t011)--(h010);
\draw[aThreeOne] (t110)--(h010);
\draw[aThreeOne] (t020)--(h010);
\draw[aThreeTwo] (h001)--(h000);
\draw[aThreeTwo] (h100)--(h000);
\draw[aThreeTwo] (h010)--(h000);

% Left input boundary.
\draw[aThreeOne]   (t002)--(left1);
\draw[aThreeTwo]   (h001)--(left2);
% The four-valent commute 13 -> 31 is absorbed into the boundary placement.
\draw[aThreeOne]   (t101)--(left4);
\draw[aThreeThree] (h000)--(left3);
\draw[aThreeTwo]   (h100)--(left5);
\draw[aThreeOne]   (t200)--(left6);

% Right input boundary.
\draw[aThreeOne]   (t002)--(right1);
\draw[aThreeTwo]   (h001)--(right2);
\draw[aThreeOne]   (t011)--(right3);
\draw[aThreeThree] (h000)--(right4);
\draw[aThreeTwo]   (h010)--(right5);
\draw[aThreeOne]   (t020)--(right6);

% Bottom output boundary.
\draw[aThreeOne]   (t200)--(bottom1);
\draw[aThreeTwo]   (h100)--(bottom2);
\draw[aThreeOne]   (t110)--(bottom3);
\draw[aThreeThree] (h000)--(bottom4);
\draw[aThreeTwo]   (h010)--(bottom5);
\draw[aThreeOne]   (t020)--(bottom6);

% Six trivalent contraction vertices.
\node[aThreeT] at (t002) {};
\node[aThreeT] at (t101) {};
\node[aThreeT] at (t011) {};
\node[aThreeT] at (t200) {};
\node[aThreeT] at (t110) {};
\node[aThreeT] at (t020) {};

% Four six-valent braid vertices.
\node[aThreeH] at (h001) {};
\node[aThreeH] at (h100) {};
\node[aThreeH] at (h010) {};
\node[aThreeH] at (h000) {};

% Tetrahedral lattice labels.
\node[aThreeVertexLabel,above=1.25mm of t002] {$(0,0,2,0)$};
\node[aThreeVertexLabel,above left=0.80mm and 0.60mm of t101]
  {$(1,0,1,0)$};
\node[aThreeVertexLabel,right=1.20mm of t011] {$(0,1,1,0)$};
\node[aThreeVertexLabel,left=1.20mm of t200]  {$(2,0,0,0)$};
\node[aThreeVertexLabel,below=1.10mm of t110] {$(1,1,0,0)$};
\node[aThreeVertexLabel,right=1.20mm of t020] {$(0,2,0,0)$};
\node[aThreeVertexLabel,right=1.20mm of h001] {$(0,0,1,1)$};
\node[aThreeVertexLabel,left=1.20mm of h100]  {$(1,0,0,1)$};
\node[aThreeVertexLabel,right=1.20mm of h010] {$(0,1,0,1)$};
\node[aThreeVertexLabel,left=1.20mm of h000]  {$(0,0,0,2)$};

% Boundary words are read counterclockwise: up the left side, down the right
% side, and from left to right along the bottom side.
\node[aThreeBoundaryLabel,rotate=60] at (-4.70,6.45)
  {input: $\boldsymbol\Delta=121321$};
\node[aThreeBoundaryLabel,rotate=-60] at (4.70,6.45)
  {input: $\boldsymbol\Delta=121321$};
\node[aThreeBoundaryLabel] at (0,-0.45)
  {output: $\boldsymbol\Delta=121321$};
\end{tikzpicture}
